\section{Introduction}
% Vulnerability
Software vulnerabilities are weaknesses in a software system that could be exploited by attackers. This exploitation can cause substantial damage, especially for the critical systems~\citep{survey_3}.
%
To reduce manual effort in discovering vulnerabilities, researchers have invested considerable effort in investigating effective approaches for \textit{automated vulnerability detection}, leading to many techniques~\citep{linevd,linevul,ivdetect,vuldeeppeaker,sysevr,velvet,poster,vulsniper,devign,are_we_there,cheng2022path}. 
%
Recent advances in this field have resulted in a (quite) high accuracy in determining whether an entire given method/file is vulnerable or not~\cite{linevul,vuldeeppeaker,sysevr,velvet,poster}. For example, with \textit{BigVul} benchmark~\cite{bigvul}, Fu \etal~\cite{linevul} propose \textit{LineVul}, which achieves more than 90\% accuracy in vulnerability detection at the function level.
%
With the support of these state-of-the-art techniques, the next step, security analysis, that developers have to perform is investigating the detected vulnerable functions to determine the actual presence of the vulnerability. 
%
% With the support of these state-of-the-art techniques, the next step, security analysis, that developers have to perform is investigating the detected vulnerable functions to determine the actual presence of the vulnerability. 
%
However, even having vulnerable code, investigating to fix those functions could still be a non-trivial task~\cite{vul_secret,openai_fixing,transfer_learning_for_fixing}. 
%

% Meanwhile, \hl{the \textit{types of a vulnerability},} such as \hl{\textit{buffer overflow} or \textit{memory corruption}, could provide valuable information to developers in debugging the vulnerable code. Indeed, this information will tell the vulnerability principles and help developers quickly pinpoint the precise location of the vulnerability and suggest potential methods to fix the vulnerable code}~\cite{mu_vuldeeppeker}. 
% %
% %
% \hl{For example, given a detected vulnerable function in Figure}~\ref{fig:BO_example}, \hl{if a developer knows that the function has a \textit{buffer overflow} error (line 30), she could quickly locate the code statements writing/copying data to buffer and investigate the possibility those statements could overrun the buffer's boundary and overwrite adjacent memory locations. Thus, she can start security analysis with the code statements at lines 7, 19, and 30 rather than all the function's statements.
% %
% %
% Once the vulnerable statement (line 30) is found, she could apply the known techniques to prevent it, such as adding checking size before writing or correcting the amount of data to be copied. Thus, after vulnerable code is detected, identifying the types of vulnerability could help reduce the workload of developers and code auditors, which is especially valuable for large vulnerable pieces of code.}



\hl{Meanwhile, the knowledge of vulnerability types, such as \textit{buffer overflow} or \textit{memory corruption}, can offer significant insights to developers when debugging the vulnerable code. This information serves as a guide to understanding the principles behind the vulnerability, enabling developers to swiftly pinpoint the exact location of the vulnerability and propose potential approaches to rectify the flawed code}~\cite{mu_vuldeeppeker}. 
%
\hl{For instance, if a developer identifies a vulnerable function, as shown in Figure}~\ref{fig:BO_example}, \hl{and knows it contains a buffer overflow error, she can immediately focus her investigation on the code statements responsible for writing or copying data into the buffer. This allows her to assess the likelihood of these statements exceeding the buffer's boundaries and overwriting adjacent memory locations. Consequently, she can initiate security analysis by examining the code statements at lines 7, 19, and 30 instead of analyzing the entire function. Once the vulnerable statement (line 30) is identified, she can employ established techniques, such as implementing size checks before writing or correcting the amount of data to be copied, to mitigate the vulnerability. Thus, by determining the vulnerability types after detecting vulnerable code, developers and code auditors can significantly reduce their workload, which is especially beneficial when dealing with large sections of vulnerable code.}


\begin{figure}
    \centering
    \lstinputlisting[language=C]{imgs/example.m}
        \caption{\hl{A \textit{buffer overflow} in Linux Kernel, CVE-2016-6156}}
        \label{fig:BO_example}
\end{figure}


Despite the importance of \hl{vulnerability type identification (VTI) in debugging vulnerable code after being detected}, the problem has not received the deserved attention. In this work, we make the first step to explore the problem of VTI by using the existing techniques in software engineering (SE) and natural language processing (NLP). Particularly, we model VTI as the multi-label text classification task in NLP~\cite{multilabel_cls,text_multi_label,brm}. \hl{This is reasonable because each vulnerable function $f$ could be considered as a document $d$, and the type set of the vulnerabilities in $f$ could be the label set of $d$}. 

The NLP community recently witnessed a dramatic paradigm shift towards the ``\textit{pre-training + fine-tuning}'' framework. Deep pre-trained models, e.g., BERT~\cite{bert}, induce powerful embeddings that can be rapidly fine-tuned on many downstream NLP problems by adding a task-specific lightweight linear layer on top of the transformer models. BERT-like models (e.g., XLNet~\cite{xlnet} and RoBERTa~\cite{roberta}) have led to state-of-the-art performance on many NLP tasks, such as part-of-speech tagging or text classification.
%
For the SE, pre-trained models have recently shown to be highly effective in many classification tasks such as bug detection, clone detection, and vulnerability detection~\cite{zhang2019novel,embedding_emse22,embedding_icse19,assessing_codebert, linevd,linevul}. This naturally raises a question: \textit{How are these pre-trained models effective for the VTI task?}

% Additionally, we suspect that given a vulnerable function $f$ and a set of vulnerability types $T$, the types of the vulnerability in $f$, $S \subseteq T$ could be determined by the appearance of code tokens in $f$. 
Additionally, we suspect that given a set of all possible vulnerability types $T$ and \hl{a vulnerable function $f$, the types of the vulnerability in $f$}, $S \subseteq T$ could be determined by the appearance of code tokens in $f$. 
%
For example, \textit{directory traversal} vulnerabilities, which allow attackers to access files/directories stored outside the web root directory, usually contain certain code (sub)tokens such as \texttt{file}, \texttt{base\_path}, or \texttt{directory}. Another example is that \textit{buffer overflow} errors usually associate with \texttt{buffer} (\texttt{buf}) or \texttt{copy} (\texttt{cpy}).
%
Meanwhile, the importance or relevance of these (sub)tokens in determining the types of vulnerabilities could be captured well by the old-fashioned code representations such as TF-IDF features. Another natural question is: \textit{How does a simple classification model with an old-fashioned code representation such as TF-IDF work for the VTI task compared to pre-trained models?}



In this paper, we conduct experiments to evaluate the performance of the state-of-the-art methods in natural language processing (NLP) and software engineering (SE) for the VTI task. We select two pre-trained models, Word2vec~\cite{word2vec_1} and CodeBERT~\cite{codebert}, which are the representative non-contextual and contextual code embeddings~\cite{ese_empirical}. We pick Word2vec because this model has become one of the most popular pre-trained models for code due to its efficiency~\cite{zhang2019novel,embedding_emse22}. Meanwhile, the reasons for our selection of CodeBERT are the model's reputation and its strong improvements in many SE tasks~\cite{assessing_codebert}. We also compare these state-of-the-art models against a simple multi-label classification model, Binary Relevance (BR), with old-fashion TF-IDF features. With the BR classifier and TF-IDF, this approach is considered as the baseline of multi-label text classification in NLP~\cite{brm} (so-called \base).

%initial results
% To evaluate VTI approaches, we used BigVul benchmark~\cite{bigvul}, which is widely used for evaluating vulnerability detection techniques~\cite{linevd,linevul,ivdetect}.
% To evaluate VTI approaches, we used BigVul benchmark~\cite{bigvul}, which contains 10K functions with vulnerabilities in 13 types.
%
% This dataset contains more than 10K functions with vulnerabilities in 13 types. 
%
Our experiments on BigVul benchmark~\cite{bigvul} show a \textit{surprising result} that \textit{the advanced pre-trained models just slightly improve the VTI performance of the baseline approach by only less than 7\% in classification accuracy}. Meanwhile, the deep pre-trained models require GPUs and cost up to 40X and 5X more time in training and predicting compared to the baseline approach. These show that for the VTI task, TF-IDF and classical Binary Relevance could capture well the features to distinguish vulnerability types and \textit{efficiently} achieve performance very competitive with the deep pre-trained models in VTI.


In this work, we also introduce a simple technique to improve the VTI performance of \base. Our idea is that certain code (sub)tokens are more likely to appear/not appear in the vulnerabilities of a type than the others. 
%
% These \textit{distinguishing (sub)tokens} could be identified beforehand (before predicting) by statistically analyzing previously known vulnerable code. 
%
%
% Particularly, we analyze the syntactic code elements (e.g., \textit{function calls}, \textit{assignments}, or \textit{control structures}), which could provide useful information to determine vulnerabilities' types.
%
These \textit{distinguishing tokens} are identified beforehand (before predicting) by statistically analyzing the syntactic code elements crucial for VTI (e.g., \textit{function calls}, \textit{assignments}, or \textit{control structures}) in previously known vulnerable code.
% such as the control structures, assignments, and function calls
As an \textit{independent component}, these distinguishing tokens are used to refine the predictions produced by \base. 
%
The intuition is that if a function is predicted to contain a vulnerability \textit{not} of type $t$, yet actually contains the distinguishing tokens of $t$ which are prevalent in the cases of $t$ but not the other types; then the prediction will be refined to the vulnerability of $t$.

Our experimental results show that the combination of \base and our component significantly improves the VTI performance of \base and outperforms the advanced pre-trained VTI models. Meanwhile, the predicting time slightly increases, and identifying distinguishing tokens does not increase the overall training and preparing time. These results indicate that \textit{\base combining with a very lightweight component could improve the VTI performance of advanced pre-trained models while retaining a very high efficiency}. We also show that our technique could effectively improve the other VTI approaches when applied on top of them up to 92.8\% F1-score.


 In brief, this paper makes the following contributions:
\begin{enumerate}
    \item An exploratory study on the performance of both traditional and advanced VTI techniques.
    \item Surprising experimental results showing that a simple model with an old-fashioned feature extraction could achieve a very competitive performance with the state-of-the-art approaches.
    \item A lightweight but effective technique improving the performance of the existing approaches.
\end{enumerate}
%
The detailed implementation and dataset can be found at: \textit{\url{https://github.com/sonnguyenvnu/VIT-Project}}.


The rest of this paper is organized as follows. Section~\ref {sec:problem} states the problem of vulnerability type identification (VTI) with a benchmark and evaluation metrics for VTI approaches. Then, several VTI approaches modeling the VTI task as the multi-label classification task and our surprising results are introduced in Section~\ref{sec:approaches}. 
% Our surprising result on the VTI performance of these approaches is shown in Section~\ref{sec:results}.
After that, Section~\ref{sec:our_approach} introduces the design of our lightweight independent technique and its effectiveness in improving \base as well as the other approaches. 
%
Some threats to validity are discussed in Section~\ref{sec:threats}.
%
Section~\ref{sec:related_work} provides the related work. Finally, Section~\ref{sec:conclusion} concludes this paper.